% Source: https://artofproblemsolving.com/community/c3h2888241_bi_parallel_to_angle_bisector
% Stable identity: aops-topic-2888241-q01
The homothety with centre the centroid $G$ and ratio $-\dfrac12$ maps $A,B,C$ to $D,E,F$ respectively. (For example, $B\mapsto E$ because $G$ divides the median $BE$ in the ratio $2:1$.) So it maps $\angle ABC$ onto $\angle DEF$.
\begin{asy}
usepackage("amsmath");
size(3.5cm);
defaultpen(fontsize(10pt));
pair A=(1,4),B=(0,0),C=(6,0);pair D=(B+C)/2,E=(A+C)/2,F=(A+B)/2;real a=abs(B-C),bb=abs(C-A),cc=abs(A-B);pair I=(a*A+bb*B+cc*C)/(a+bb+cc);
draw(A--B--C--cycle);draw(D--E--F--cycle,gray);draw(B--I,red);pair u=unit(unit(D-E)+unit(F-E));draw(E--E+1.6*u,red+dashed);dot(I);
label("$A$",A,N);label("$B$",B,SW);label("$C$",C,SE);label("$D$",D,S);label("$E$",E,NE);label("$F$",F,W);label("$I$",I,N);
\end{asy}

A homothety maps angle bisectors to angle bisectors and each line to a parallel line. So the bisector of $\angle DEF$ is parallel to the bisector of $\angle ABC$, which is $BI$. $\blacksquare$
